\documentclass[11pt,a4paper]{article}

\usepackage[french]{babel}
\usepackage{amsmath,amssymb,bm}
\usepackage{physics}
\usepackage{geometry}

\geometry{margin=2.5cm}

\title{Énergie d'échange et anisotropie dans un cylindre infini (mode de curling)}
\author{}
\date{}

\begin{document}

\maketitle

%%%%%%%%%%%%%%%%%%%%%%%%%%%%%%%%%%%%%%%%%%%%%%%%%%%%%%%%%%%%%%%

\section{Aimantation}

On considère un cylindre ferromagnétique infini de rayon $R$.
Les coordonnées cylindriques sont $(r,\varphi,z)$.

L'aimantation réduite est supposée de la forme :

\[
\boxed{
\mathbf m(r)
=
\sin\omega(r)\,\mathbf e_\varphi
+
\cos\omega(r)\,\mathbf e_z
}
\]

avec $\omega(r)$ l'angle par rapport à $Oz$.

On a bien $|\mathbf m|=1$.

%%%%%%%%%%%%%%%%%%%%%%%%%%%%%%%%%%%%%%%%%%%%%%%%%%%%%%%%%%%%%%%

\section{Gradient tensoriel}

Le gradient tensoriel en coordonnées cylindriques est :

\[
\nabla\mathbf m
=
\frac{\partial\mathbf m}{\partial r}\otimes\mathbf e_r
+
\frac1r\frac{\partial\mathbf m}{\partial\varphi}\otimes\mathbf e_\varphi
+
\frac{\partial\mathbf m}{\partial z}\otimes\mathbf e_z.
\]

Les dérivées des vecteurs de base sont :

\[
\frac{\partial\mathbf e_r}{\partial\varphi}=\mathbf e_\varphi,\quad
\frac{\partial\mathbf e_\varphi}{\partial\varphi}=-\mathbf e_r.
\]

%%%%%%%%%%%%%%%%%%%%%%%%%%%%%%%%%%%%%%%%%%%%%%%%%%%%%%%%%%%%%%%

\subsection{Dérivée radiale}

\[
\frac{\partial\mathbf m}{\partial r}
=
\omega'(r)\cos\omega\,\mathbf e_\varphi
-
\omega'(r)\sin\omega\,\mathbf e_z
\]

\[
\left|\frac{\partial\mathbf m}{\partial r}\right|^2=(\omega')^2.
\]

%%%%%%%%%%%%%%%%%%%%%%%%%%%%%%%%%%%%%%%%%%%%%%%%%%%%%%%%%%%%%%%

\subsection{Dérivée azimutale}

\[
\frac{\partial\mathbf m}{\partial\varphi}
=
-\sin\omega\,\mathbf e_r
\]

\[
\left|\frac1r\frac{\partial\mathbf m}{\partial\varphi}\right|^2
=
\frac{\sin^2\omega}{r^2}.
\]

%%%%%%%%%%%%%%%%%%%%%%%%%%%%%%%%%%%%%%%%%%%%%%%%%%%%%%%%%%%%%%%

\subsection{Gradient et norme}

\[
\boxed{
\nabla\mathbf m:\nabla\mathbf m
=
(\omega')^2+\frac{\sin^2\omega}{r^2}
}
\]

%%%%%%%%%%%%%%%%%%%%%%%%%%%%%%%%%%%%%%%%%%%%%%%%%%%%%%%%%%%%%%%

\section{Énergie d'échange}

\[
w_{\rm ex}
=
A\left[
(\omega')^2+\frac{\sin^2\omega}{r^2}
\right]
\]

\[
\boxed{
E_{\rm ex}
=
2\pi A
\int_0^R
\left[
r(\omega')^2+\frac{\sin^2\omega}{r}
\right]dr
}
\]

%%%%%%%%%%%%%%%%%%%%%%%%%%%%%%%%%%%%%%%%%%%%%%%%%%%%%%%%%%%%%%%

\section{Ajout d'une anisotropie uniaxiale}

On ajoute une anisotropie uniaxiale d'axe $Oz$ et de constante $K_u$.

L'énergie volumique est :

\[
w_{\rm an}
=
K_u(1-m_z^2)
\]

or

\[
m_z=\cos\omega
\]

donc

\[
\boxed{
w_{\rm an}=K_u\sin^2\omega
}
\]

%%%%%%%%%%%%%%%%%%%%%%%%%%%%%%%%%%%%%%%%%%%%%%%%%%%%%%%%%%%%%%%

\section{Énergie d'anisotropie}

\[
\boxed{
E_{\rm an}
=
2\pi K_u
\int_0^R
\sin^2\omega(r)\,r\,dr
}
\]

%%%%%%%%%%%%%%%%%%%%%%%%%%%%%%%%%%%%%%%%%%%%%%%%%%%%%%%%%%%%%%%

\section{Énergie totale}

La densité totale est :

\[
w=
A\left[
(\omega')^2+\frac{\sin^2\omega}{r^2}
\right]
+
K_u\sin^2\omega
\]

Donc :

\[
\boxed{
E
=
2\pi
\int_0^R
\left[
Ar(\omega')^2
+
\left(\frac{A}{r}+K_u r\right)\sin^2\omega
\right]dr
}
\]

%%%%%%%%%%%%%%%%%%%%%%%%%%%%%%%%%%%%%%%%%%%%%%%%%%%%%%%%%%%%%%%

\section{Équation d'Euler-Lagrange}

Le lagrangien radial est :

\[
L=
Ar(\omega')^2
+
\left(\frac{A}{r}+K_u r\right)\sin^2\omega
\]

L'équation d'Euler-Lagrange donne :

\[
\frac{d}{dr}(2Ar\omega')
-
2\left(\frac{A}{r}+K_u r\right)\sin\omega\cos\omega=0
\]

soit :

\[
\boxed{
\omega''
+\frac1r\omega'
-
\left(\frac1{r^2}+\frac{K_u}{A}\right)
\sin\omega\cos\omega
=0
}
\]

En introduisant la variable adimensionnelle  $\rho = \frac{r}{R}$, l'équation différentielle devient:
\[
\omega_{\rho\rho}
+\frac{1}{\rho}\omega_\rho
-\left(\frac{1}{2 \rho^2} + \kappa\right)\sin(2\omega)
= 0
\]

avec \[
\kappa = \frac{K_u R^2}{2A}
\]

Les conditions aux limites s'écrivent:
\[
\omega(\varepsilon)=0,
\qquad
\omega_\rho(1)=0.
\]

et en notant $\epsilon =  \frac{r_0}{R}$.

%%%%%%%%%%%%%%%%%%%%%%%%%%%%%%%%%%%%%%%%%%%%%%%%%%%%%%%%%%%%%%%

\section{Remarques physiques}

\begin{itemize}

\item Le terme $\frac{1}{r^2}$ provient de la géométrie cylindrique (rotation de la base).

\item Le terme $\frac{K_u}{A}$ joue le rôle d'une longueur d'échelle d'anisotropie.

\item L'équation ci-dessus est celle utilisée pour l'étude du mode de curling et du champ de nucléation.

\end{itemize}

\section*{Justification du champ et de l'énergie démagnétisante}

Le champ démagnétisant $\vec{H}_d$ est dicté par les équations de la magnétostatique. En l'absence de courants libres, il dérive d'un potentiel scalaire magnétique et sa divergence est directement reliée à la divergence de l'aimantation :
\begin{equation}
    \vec{\nabla} \cdot \vec{H}_d = -M_s\; \vec{\nabla} \cdot \vec{m} = \rho_m
\end{equation}
où $\rho_m$ représente la densité volumique de charges magnétiques. En coordonnées cylindriques, compte tenu des invariances de notre système par translation selon $Oz$ et par rotation autour de ce même axe, l'aimantation s'écrit $\vec{m} = m_\rho(\rho)\vec{e}_\rho + m_\phi(\rho)\vec{e}_\phi + m_z(\rho)\vec{e}_z$. Sa divergence se réduit alors strictement à sa contribution radiale :
\begin{equation}
    \vec{\nabla} \cdot \vec{m} = \frac{1}{\rho} \frac{\text{d}}{\text{d}\rho} \left(\rho \, m_\rho(\rho)\right) \implies \rho_m(\rho) = -M_s\frac{1}{\rho} \frac{\text{d}}{\text{d}\rho} \left(\rho \, m_\rho(\rho)\right)
\end{equation}
Par application du théorème de Gauss magnétique à un cylindre de rayon $\rho$ coaxial au système, le champ démagnétisant, purement radial par symétrie ($\vec{H}_d = H_d(\rho)\vec{e}_\rho$), est déterminé par la charge interne cumulée. Sachant qu'à l'axe ($\rho=0$) l'aimantation pointe selon $Oz$ (ce qui implique $m_\rho(0)=0$), l'intégration de la densité volumique donne de manière univoque :
\begin{equation}
    \vec{H}_d(\rho) = -M_s \; m_\rho(\rho)\vec{e}_\rho
\end{equation}
La densité locale d'énergie démagnétisante par unité de volume s'exprime par le produit scalaire standard :
\begin{equation}
    E_{\text{dém}} = -\frac{1}{2}\mu_0 \vec{H}_d \cdot \vec{m} = \frac{1}{2}\mu_0 M_s^2 \; m_\rho(\rho)^2
\end{equation}

\section*{Choix optimal pour réduire l'énergie}

En exprimant les composantes de l'aimantation de norme constante $M_s$ en fonction de l'angle polaire $\omega$ (avec $Oz$) et de l'angle azimutal local $\phi$ (avec $\vec{e}_\rho$), la composante radiale devient $m_\rho = M_s \sin\omega \cos\phi$. L'énergie micromagnétique totale par unité de volume, incluant l'échange et l'anisotropie uniaxiale $K_u$, s'écrit donc :
\begin{equation}
    E = A \left[ \left(\frac{\text{d}\omega}{\text{d}\rho}\right)^2 + \sin^2\omega \left(\frac{\text{d}\phi}{\text{d}\rho}\right)^2 + \frac{\sin^2\omega}{\rho^2} \right] + K_u \sin^2\omega + \frac{1}{2}\mu_0 M_s^2 \sin^2\omega \cos^2\phi
\end{equation}

Une analyse de cette fonctionnelle d'énergie permet de dégager les propriétés suivantes pour l'état d'équilibre :
\begin{itemize}
    \item \textbf{Verrouillage de l'angle azimutal :} Le terme d'énergie démagnétisante est proportionnel à $\cos^2\phi$. Étant donné que le coefficient de l'énergie dipolaire est strictement positif ($\mu_0 M_s^2 > 0$), le système minimise rigoureusement son énergie globale en annulant ce terme. Cela impose la condition \textbf{$\phi = \pi/2$} en tout point du cylindre.
    \item \textbf{Configuration en vortex pur (mode \textit{curling}) :} Physiquement, la condition $\phi = \pi/2$ signifie que la composante radiale de l'aimantation s'annule ($m_\rho = 0$). L'aimantation s'enroule de manière purement azimutale et axiale. D'après l'analyse magnétostatique préalable, ce choix annule simultanément le champ $\vec{H}_d$, les charges volumiques ($\rho_m = 0$) ainsi que les charges de surface en $\rho = R$ ($\sigma_m = \vec{m} \cdot \vec{e}_\rho = 0$). L'énergie démagnétisante est donc rigoureusement nulle à l'équilibre.
    \item \textbf{Équation d'Euler-Lagrange réduite :} En substituant $\phi = \pi/2$ et $\frac{\text{d}\phi}{\text{d}\rho} = 0$, la compétition énergétique se résume à l'équilibre entre la raideur d'échange et l'anisotropie. L'équation d'Euler-Lagrange minimisant l'énergie par rapport à l'angle polaire $\omega(\rho)$ devient :
    \begin{equation}
        \frac{\text{d}^2\omega}{\text{d}\rho^2} + \frac{1}{\rho}\frac{\text{d}\omega}{\text{d}\rho} - \frac{\sin(2\omega)}{2\rho^2} - \frac{K_u}{2A}\sin(2\omega) = 0
    \end{equation}
\end{itemize}

En introduisant la coordonnée radiale adimensionnelle $\tilde{\rho} = \rho/R$ et le paramètre d'anisotropie réduit $\kappa = \frac{K_u R^2}{2A}$, l'équation différentielle prend la forme finale implémentée dans le code de résolution numérique :
\begin{equation}
    \frac{\text{d}^2\omega}{\text{d}\tilde{\rho}^2} = -\frac{1}{\tilde{\rho}}\frac{\text{d}\omega}{\text{d}\tilde{\rho}} + \left(\frac{1}{2\tilde{\rho}^2} + \kappa\right)\sin(2\omega)
\end{equation}
Puisque l'anisotropie considérée est de type plan facile ($K_u < 0$), le paramètre $\kappa$ est négatif. Ce terme pousse l'aimantation à s'écarter de l'axe longitudinal pour s'allonger dans le plan ($\omega \to \pi/2$) à mesure que l'on s'éloigne du cœur du cylindre.

\end{document}